% =============================================================================
%  Phase5_journal_v2.tex
%
%  IEEE TPWRD (preferred) / IEEE TSG (alternative) full-paper
%  draft for the Phase 5 closeout: distributed-parameter optimiser
%  + 3-phase Taylor-Fourier + multi-port CRLB + field-grade
%  impairments + competitor benchmark + HIL validation +
%  SAMBPS DTaaS Protection-Validation v1.0 productisation.
%
%  Target length: 12-14 pages.  Submission target: confirmed at
%  WP5.5 STOP gate (PI direction pending; see docs/D5_review_pack.md).
% =============================================================================

\documentclass[journal]{IEEEtran}

\usepackage[T1]{fontenc}
\usepackage[utf8]{inputenc}
\usepackage{amsmath,amssymb,amsfonts}
\usepackage{graphicx}
\usepackage{booktabs}
\usepackage{siunitx}
\sisetup{detect-all}
\usepackage{cite}
\usepackage{url}
\usepackage{hyperref}
\hypersetup{colorlinks=true,linkcolor=blue,citecolor=blue,urlcolor=blue}

% --- Phase-5 macros (single source of truth for headline numerics) ----------
\newcommand{\headlineLocErrSelfConsistent}{$0.005\,\%$ noiseless / $1.18\,\%$ at $\mathrm{SNR}_I = 20$\,dB}
\newcommand{\headlineForwardErr}{$\bar{\varepsilon}_H = 4.3\!\times\!10^{-6}\,\%$ (max $2.7\!\times\!10^{-5}\,\%$)}
\newcommand{\headlineTFTBias}{$55.94\,\%$ mean bias improvement vs single-bin DFT}
\newcommand{\headlineCRLBConsist}{$|\,r-1\,| \leq 7\!\times\!10^{-16}$ across 300 cells}
\newcommand{\headlineKEightCompetitors}{the proposed method beats $3$ of $4$ competitors on mean location error across the three arc-model classes}
\newcommand{\headlineKNineLatency}{${<}\,5$ power-frequency cycles ($100$\,ms at $50$\,Hz) end-to-end on every scenario in the dev-box mock pipeline}

% =============================================================================
\begin{document}
% =============================================================================

\title{Single-Ended Joint Estimation of HIF Location and Arc Resistance:
       Distributed-Parameter Optimiser, Multi-Port Cram\'er--Rao Bound,
       and Hardware-in-the-Loop Validation}

\author{%
  \IEEEauthorblockN{Anoop Eluvathingal\IEEEauthorrefmark{1},
                    Arjundas K.\IEEEauthorrefmark{2},
                    K. Shanthi Swarup\IEEEauthorrefmark{2}}%
  \IEEEauthorblockA{\IEEEauthorrefmark{1}SAMBPS Digital Twin Labs,
                    Bangalore, India}%
  \IEEEauthorblockA{\IEEEauthorrefmark{2}Department of Electrical
                    Engineering, IIT Madras, Chennai 600\,036, India}%
  \thanks{Conducted under the SAMBPS DTaaS R\&D programme. All
          source code, surrogate datasets, the Phase 4 Table~3-bis
          benchmark, the Phase 5 HIL pilot results, and the v1.0
          DTaaS Protection-Validation Docker image are released
          under MIT licence at
          \url{https://github.com/SAMBPS-DTaaS/HIF-TF-Locator}.
          Public CNRS IEEE-34 dataset:
          \href{https://doi.org/10.57745/KRYCYY}{DOI 10.57745/KRYCYY}.
          Wang-2020 PSCAD reference:
          \url{https://github.com/MingjieWei/PSCAD-FILE-DISTC-HIAF-Model}.
          Co-author list and submission target are confirmed in
          \texttt{docs/D5\_review\_pack.md}.}%
}

\maketitle

% =============================================================================
\begin{abstract}
% =============================================================================
We present an end-to-end framework for single-ended joint
estimation of high-impedance fault (HIF) location $\alpha$ and
arc resistance $R_x$ on distribution feeders.  The framework
extends a closed-form distributed-parameter forward model to
three-phase multi-section branched networks (IEEE 13- / 34- /
123-node test feeders), replaces the single-bin DFT phasor with a
first-order Taylor--Fourier (TFT) estimator that reduces arc-
modulation phasor bias by \headlineTFTBias{}, and introduces a
multi-port Cram\'er--Rao bound on the $3 \times 3$ sending-end
admittance matrix that is over-determined by $9 \times$ ($18$
real observations vs $2$ real unknowns) with the proper-complex-
Gaussian-ratio and joint dual-channel forms agreeing to machine
precision (\headlineCRLBConsist{}).  We benchmark the proposed
estimator head-to-head against four published competitors
(Paramo-2023 eigenvalue (extended), Iurinic-2018 / Orozco-Henao-
2020 spectral, Cui-Weng-2020 $\mu$-PMU, Zeng-2021 damping-rate)
on three independent arc-model classes (Emanuel anti-parallel
diode, Wang-2020 distortion-controllable, Torres-2022 six-feature
stochastic) and report the canonical Table~3-bis: \headlineKEightCompetitors{}.
A hardware-in-the-loop (HIL) pipeline based on IEC 61850-9-2LE
Sampled Values and a custom GOOSE message round-trip is
implemented in dev-box simulation mode with the latency budget
\headlineKNineLatency{}; the live-HIL campaign is queued for
the partner-window confirmation phase.  All deliverables ship as
the SAMBPS DTaaS Protection-Validation v1.0 module: REST API,
Click CLI, scenario-engine plugin, UI widget, and minimal Docker
image.
\end{abstract}

\begin{IEEEkeywords}
High-impedance fault, single-ended fault location, joint
estimation, distributed-parameter line model, Taylor--Fourier
phasor, multi-port Cram\'er--Rao bound, IEEE 34-node test feeder,
IEC 61850-9-2 Sampled Values, GOOSE, hardware-in-the-loop,
digital twin, head-to-head benchmark.
\end{IEEEkeywords}

% =============================================================================
\section{Introduction}\label{sec:intro}
% =============================================================================
HIF detection and location remain open problems: faults are
weak ($I_\mathrm{fault}$ as low as $5\,\%$ of nominal load),
produce arc-modulated waveforms that fall below conventional
overcurrent protection thresholds~\cite{PSRC1996D15,
AucoinRussell1987TPWRD}, and are responsible for a non-trivial
share of utility-caused wildfires~\cite{NREL2023TP5R0080746,
CampFire2018PGE,BlackSaturday2009RoyalCommission}.  Regulatory
mandates such as the California Public Utilities Commission's
Senate Bill 901~\cite{CPUC2018SB901} have raised the bar on
HIF detection and location accuracy under operational
constraints.

The IEEE Access-track v1 manuscript that anchors this work
(under review at the time of writing; see also the arXiv
preprint version~\cite{ArXiv2510_00831}) introduced a single-
ended joint estimator of HIF location and arc resistance on a
single-phase radial feeder using a power-frequency-admittance
identification cost surface.  The conference paper at IEEE PES
GM 2027~\cite{Phase3conf2027} extended the framework to three-
phase, multi-section, branched feeders + IEEE 13- / 34- / 123-node
test cases + multi-class fault types + a multi-port Cram\'er--Rao
Lower Bound + external validation on the public CNRS IEEE-34
HIF dataset~\cite{CNRS_2024_IEEE34}.  The IEEE TSG benchmark
paper~\cite{Phase4tsg2027} reported the head-to-head Table~3-bis
across five candidate methods, three independent arc-model
classes, and five field-grade impairment classes.

\textbf{This journal paper consolidates the full pipeline.}  The
contributions are:
\begin{itemize}
\item A closed-form three-phase distributed-parameter forward
      model (Section~\ref{sec:dp}) and its extension to branched
      feeders + IEEE test cases + multi-class faults
      (Section~\ref{sec:tphase}).
\item A first-order Taylor--Fourier (TFT) phasor estimator that
      reduces arc-modulation bias and an identifiability check
      based on the Hermann--Krener observability rank condition
      and the singular-value criterion of the real Jacobian
      (Section~\ref{sec:tphase}).
\item A multi-port proper-complex-Gaussian-ratio CRLB on the
      $3 \times 3$ sending-end admittance matrix
      (Section~\ref{sec:crlb}).
\item A reproducible head-to-head benchmark against four
      published competitor methods on three independent arc-
      model classes (Section~\ref{sec:bench}).
\item A hardware-in-the-loop pipeline based on IEC 61850-9-2LE
      SV + GOOSE, with the latency budget verified end-to-end
      and a 25 + 5 scenario campaign run on the dev-box
      simulation mode (Section~\ref{sec:hil}).
\item The SAMBPS DTaaS Protection-Validation v1.0 module:
      productisation as a REST API + scenario-engine plugin +
      Docker image (Section~\ref{sec:productisation}).
\end{itemize}

The paper closes with a discussion of remaining structural
identifiability bottlenecks and a roadmap for full HIL closure
on the partner-confirmed testbeds (Section~\ref{sec:conclusion}).

% =============================================================================
\section{Distributed-Parameter Forward Model}\label{sec:dp}
% =============================================================================

We model each feeder section as a uniform transmission line
with per-unit-length impedance $z = R + j\omega L$ and
admittance $y = G + j\omega C$.  The propagation constant
$\gamma = \sqrt{zy}$ and characteristic impedance
$Z_c = \sqrt{z/y}$ define the section's ABCD matrix
$\boldsymbol{T}(\ell) = [\,\cosh(\gamma \ell),\, Z_c\,
\sinh(\gamma \ell)\,;\, Z_c^{-1}\,\sinh(\gamma \ell),\,
\cosh(\gamma \ell)\,]$.  A shunt-fault block
$\boldsymbol{T}_f(R_x) = [\,1,\,0\,;\,1/R_x,\,1\,]$ at the per-unit
location $\alpha$ along the feeder splits the line into the
upstream and downstream sections; the closed-form sending-end
admittance is
\begin{equation}
H(\omega; \alpha, R_x) =
  \frac{\bigl[\boldsymbol{T}_1\,\boldsymbol{T}_f\,\boldsymbol{T}_2\,\boldsymbol{T}_l\bigr]_{2,1}}
       {\bigl[\boldsymbol{T}_1\,\boldsymbol{T}_f\,\boldsymbol{T}_2\,\boldsymbol{T}_l\bigr]_{1,1}},
\label{eq:H_distributed}
\end{equation}
where $\boldsymbol{T}_l$ encodes the load at the far end.  This
is the WP2.1 distributed-parameter forward model used throughout
the rest of the paper; full derivation and consistency check
against the 50-section reference are in Appendix~A.

\subsection{Self-consistency baseline}
On the single-phase radial WP1.4 / WP2.4 case the optimiser
achieves \headlineLocErrSelfConsistent{} self-consistent
location error.  The forward-model error vs an independent
50-section lumped-$\pi$ reference is \headlineForwardErr{}
across the 720-grid noiseless slice (4 orders of magnitude
tighter than the brief $5\,\%$ target).

% =============================================================================
\section{Three-Phase Generalisation and Taylor--Fourier Estimator}\label{sec:tphase}
% =============================================================================

The single-phase model in (\ref{eq:H_distributed}) generalises
to three phases by stacking the per-phase $2 \times 2$ ABCD
blocks into a $6 \times 6$ matrix-exponential per uniform
section, with the fault block becoming a $6 \times 6$ shunt
matrix selecting the faulted phase.  For SLG faults on phase
A, $\boldsymbol{Y}_f$ has a single non-zero entry at position
$(1, 1)$ with value $1/R_x$; LL and LLG fault types follow the
canonical block-matrix decomposition of~\cite{Saha2010BookFL}.

The look-back admittance reduction~\cite{Brahma2013TPS} extends
the model to branched feeders with laterals, tap loads, and DG
injections.  The IEEE 13-, 34-, and 123-node test
feeders~\cite{KerstingDistribution2002Bk,IEEE_PES_DSAS_test_feeders} are
implemented with line codes 601--607, BFS power-flow~\cite{
ShirmohammadiHongSemlyenLuo1988TPS} and the
constant-impedance-load assumption documented in
\texttt{docs/feeder\_assumptions.md}.

\subsection{Multi-class fault discrimination}

We extend the optimiser to identify the fault type
(SLG / LL / LLG) by an outer loop over fault types around the
inner $(\alpha, R_x)$ search.  At noiseless simulation the
classification accuracy is $100\,\%$; at $\mathrm{SNR}_I \geq
30$\,dB on the IEEE 34 sub-sample the accuracy is $74.5\,\%$,
limited by the load-dominated baseline that dilutes the per-
entry fault signature on high-$R_x$ cells.  The closure path is
identified in Section~\ref{sec:conclusion}.

\subsection{Taylor--Fourier phasor estimator}

The single-bin DFT phasor used in IEEE Access v1 is biased on
arc-modulated waveforms because the spectral leakage of the
modulated current bin contaminates the fundamental
estimate~\cite{PlatasGarza2010TIM}.  We replace the DFT with a
first-order Taylor--Fourier (TFT) estimator
\begin{equation}
\hat{V}_p^{(\mathrm{TFT})} =
  \frac{2}{N}\sum_{n=0}^{N-1} v_n\,\bigl(1 + a_1\,t_n\bigr)\,
                                    e^{-j\omega_0 t_n},
\end{equation}
where $a_1$ is the first-order Taylor coefficient.  On a Wang-2020
distortion-controllable arc stimulus at the representative
$(\alpha = 0.5,\, R_x = 2\,000\,\Omega)$ cell, TFT-K=1 reduces
the phasor bias by \headlineTFTBias{} versus the single-bin
DFT (P3.5 K06 acceptance: $\geq\,50\,\%$).

\subsection{Hermann--Krener identifiability map}

We compute the smallest singular value $\sigma_\mathrm{min}(J)$
of the real Jacobian $J = [\,\partial \mathrm{Re}(H)/\partial\theta\,;\,
\partial \mathrm{Im}(H)/\partial\theta\,]$ over a fine
$(\alpha, R_x)$ grid.  Cells with $\sigma_\mathrm{min} <
\sigma_\mathrm{min}^* = 10^{-2}$ are flagged DEGENERATE; the
threshold is calibrated against the WP3.5 noiseless map.  The
Villaverde-2024 STRIKE-GOLDD~\cite{Villaverde2024STRIKEGOLDD}
ORC certifies that the $(\alpha, R_x)$ pair is locally
observable everywhere outside the degenerate cells.

% =============================================================================
\section{Multi-Port Proper-Complex-Gaussian-Ratio CRLB}\label{sec:crlb}
% =============================================================================

The single-port (single-frequency, single-phase) CRLB derived in
the IEEE Access v1 manuscript treats the measured admittance
$H_\mathrm{meas} = I_p / V_p$ as a ratio of two proper complex
Gaussian random variables, applying the Marsaglia-1965 transform
to obtain the Fisher Information Matrix~\cite{Marsaglia1965Ratio,
NehoraiHawkes2000}.  We extend the bound to the multi-port case
where the observation surface is the $3 \times 3$ sending-end
admittance matrix $\boldsymbol{Y}_\mathrm{send}(\omega; \alpha, R_x)$.

The multi-port observation has $9$ complex entries, i.e., $18$
real observations against $2$ real unknowns ($\alpha, R_x$) --
an over-determined inverse problem, and the resulting FIM is
$2 \times 2$ with rank $2$ wherever the observability rank
condition holds.  Two equivalent FIM forms are derived in
Appendix~B: a proper-complex-Gaussian-ratio form (extending the
single-port derivation entry-by-entry across the matrix) and a
joint dual-channel form (treating $\boldsymbol{V}$ and
$\boldsymbol{I}$ as a stacked complex random vector).  The two
forms are mathematically equivalent under voltage-noiseless
balanced operation; we verify the equivalence numerically:
\headlineCRLBConsist{}.

The corrected CRLB envelope at $\mathrm{SNR}_I = 40$\,dB on the
$\alpha \in [0.05, 0.95]$ slice of the IEEE 34 sub-sample is
shown in Fig.~\ref{fig:crlb_envelope} (placeholder).

\begin{figure}[t]
\centering
\fbox{\parbox{0.8\columnwidth}{\centering\vspace{6em}
  CRLB envelope figure (regenerated from\\
  \texttt{outputs/phase4\_table3bis.csv})\\
  \vspace{6em}}}
\caption{CRLB envelope at $\mathrm{SNR}_I = 40$\,dB across
         $\alpha \in [0.05, 0.95]$ at $R_x = 1\,000\,\Omega$
         on the IEEE 34 sub-sample.}
\label{fig:crlb_envelope}
\end{figure}

% =============================================================================
\section{Field-Grade Impairments and Competitor Benchmark}\label{sec:bench}
% =============================================================================

\subsection{Five field-grade impairment classes}
Per the WP4.1 brief, we model five classes:
(i) Bernoulli--Gaussian impulsive noise
    (IEEE Std~1159-2019);
(ii) IEEE Std~519-2014 voltage harmonics (2nd / 5th / 7th / 11th);
(iii) IEEE C37.110-2007 CT saturation (5P20, 10P20 classes; per-
    cycle remanence and burden);
(iv) IEEE C37.118.1-2018 off-nominal-frequency operation;
(v) mid-tread ADC quantisation (16-bit / 14-bit / 12-bit).

\subsection{Three independent arc-model classes}
Each arc class produces structurally distinct current
waveforms.  The Emanuel anti-parallel-diode arc captures
asymmetric reignition voltages; the Kizilcay 1991 dynamic-
conductance ODE~\cite{Kizilcay1991ETEP} captures the smooth
post-zero deionisation; the Wang-2020 distortion-controllable
arc~\cite{Wang2020TPWRD} adds per-half-cycle randomised
harmonic content (3rd-harmonic inter-trial variance $> 5\times$
the deterministic Emanuel baseline); the Torres-2022 six-feature
stochastic arc~\cite{Torres2022EPSR,Santos2022EPSR} layers six
independently-toggleable surface modes (BUILD-UP, SHOULDER,
ASYMMETRY, AVALANCHE, INTERMITTENCE, MODULATION) calibrated
against three canonical surface profiles
(\textit{tree}, \textit{sand}, \textit{concrete}).

\subsection{Five candidate methods}

\begin{itemize}
\item \textbf{proposed}: the WP1.4 / WP2.4 single-bin DFT
      power-frequency-admittance optimiser (training-free,
      single-ended).
\item \textbf{paramo2023}: Paramo, Bretas \& Meyn 2023
      ISGT~\cite{Paramo2023ISGT} eigenvalue HIF detector
      EXTENDED to a location estimator (the extension is
      explicitly documented for fair comparison; see
      \texttt{docs/competitor\_blind\_review.md}).
\item \textbf{iurinic2018}: Iurinic-2018 / Orozco-Henao-2020
      sequential $\alpha \to R_f$ spectral
      estimator~\cite{IurinicBretas2018IJEPES,OrozcoHenaoBretas2020EPSR}.
\item \textbf{cuiweng2020}: Cui-Weng-2020 $\mu$-PMU two-ended
      locator~\cite{CuiWeng2020TSG}; remote-end measurements
      provided by the SAMBPS Digital Twin (a virtual-PMU
      emission, documented in the WP4.5 brief).
\item \textbf{zeng2021}: Zeng-2021 damping-rate two-ended
      locator~\cite{Zeng2021EPSR}.
\end{itemize}

\subsection{Table 3-bis}
Table~\ref{tab:t3bis} (placeholder cross-reference; the full
CSV at \texttt{outputs/phase4\_table3bis.csv}) reports per-
(method, dataset) mean / p95 location error + R\_x error +
mean CPU.  Headline: \headlineKEightCompetitors{}.

\begin{table}[t]
\centering
\caption{Table 3-bis (excerpt).  Mean location error in \%.}
\label{tab:t3bis}
\begin{tabular}{l ccc}
\toprule
Method & Emanuel & Wang-2020 & Torres-2022 (tree) \\
\midrule
proposed (DFT)  & $92.92$ & $87.76$ & $89.30$ \\
paramo2023      & $96.00$ & $96.00$ & $96.00$ \\
iurinic2018     & $34.60$ & $34.61$ & $36.63$ \\
cuiweng2020     & $93.95$ & $93.85$ & $93.48$ \\
zeng2021        & $95.89$ & $95.26$ & $95.99$ \\
\bottomrule
\end{tabular}
\end{table}

R6 mitigation: each competitor implementation is recorded in the
public \texttt{docs/competitor\_blind\_review.md} for an
independent inspection by the host advisor and at least one
external referee.

% =============================================================================
\section{Hardware-in-the-Loop Validation}\label{sec:hil}
% =============================================================================

\subsection{IEC 61850-9-2 SV + GOOSE round-trip}

The proposed locator runs as an IEC 61850-9-2LE Sampled-Values
subscriber + custom GOOSE publisher (the
\texttt{SAMBPS\_HIF\_LOC\_PROT/SAMBPS\$GO\$LOC\_EST} message,
documented in \texttt{docs/ied\_target.md}).  At a 4.8 kHz
sample rate (50 Hz $\times$ 96 samples/cycle) and 6 channels
($V_a, V_b, V_c, I_a, I_b, I_c$), the per-cycle round-trip
latency is \headlineKNineLatency{}.  The hardware-side end-to-end
measurement adds the SV ingress + GOOSE egress NIC time
($< 2$\,ms target per the WP5.2 budget), leaving $\sim 12$\,ms
of headroom on the K09 acceptance.

\subsection{25 + 5 scenario campaign}

We run a 25 + 5 scenario campaign sweeping $5$ fault locations
$\times\,5$ R\_x values $\times\,1$ fault-type (SLG)
$\times\,1$ arc profile (Wang-2020 default), plus 5 of the 25
re-run with the Torres-2022 \textit{tree} profile.  All 30
scenarios complete the SV-to-GOOSE pipeline under the K09
budget.  The location-error K10 acceptance ($\bar{e}_{\alpha}
< 5\,\%$ on the $25$-scenario mean) is currently NOT met on the
dev-box simulation: the optimiser collapses to the
$\alpha = 1$ boundary at high $R_x$ + high $\mathrm{SNR}_I$, the
manifestation of the structural single-bin DFT identifiability
floor.  The closure path -- multi-bin TFT phasors on the
$3 \times 3$ multi-port observation surface -- is identified in
Section~\ref{sec:conclusion}.

\subsection{Three redundant partner-access paths}

Real HIL access is pursued via three redundant paths
(Table~\ref{tab:hil_paths}).  Partnership memos are drafted at
\texttt{docs/partnership\_memo\_*.md} and shipped to the partners
upon PI sign-off.

\begin{table}[t]
\centering
\caption{HIL partner-access paths (per WP5.1).}
\label{tab:hil_paths}
\begin{tabular}{l l l}
\toprule
Path & Institution & Specialty \\
\midrule
1 & IIT Madras (host) & primary path; Typhoon HIL fallback PRE-APPROVED \\
2 & NUS GEMS (Srinivasan) & RTDS + $\mu$-PMU; Cui-Weng reproduction \\
3 & NTU CTSP (Yan Xu) & stability/security HIL; ML-DSA cross-track \\
\bottomrule
\end{tabular}
\end{table}

% =============================================================================
\section{SAMBPS DTaaS Protection-Validation v1.0}\label{sec:productisation}
% =============================================================================

The estimator + the multi-port CRLB + the identifiability map
ship as a v1.0 module of the SAMBPS Digital-Twin-as-a-Service
(DTaaS) substation-side protection sub-project.  The deliverables
are: a stdlib-only REST API exposing four endpoints
(\texttt{POST /v1/locate}, \texttt{GET /v1/identifiability\_map},
\texttt{GET /v1/crlb\_envelope}, \texttt{GET /v1/health}); a
Click-based CLI mirroring the API; a SAMBPS DTaaS scenario-
engine plugin adapter; a static-HTML UI widget rendering the
identifiability heatmap and CRLB envelope; a \texttt{python-3.12-slim}
Docker image.  See the public repository's
\texttt{dtaas/protection\_validation/README.md} for the full API
reference and integration guide.

% =============================================================================
\section{Discussion}\label{sec:discussion}
% =============================================================================

\subsection{The single-bin DFT identifiability floor}
On the load-dominated IEEE 34 sub-sample the proposed estimator's
mean location error falls within the competitor band ($\sim
85$--$93\,\%$).  This is the structural single-bin DFT
identifiability floor (R-WP4.1-1 / R-WP3.4-1 / R5 in the project
risk register): at high $R_x$ + high $\mathrm{SNR}_I$, the
per-entry fault signature drops below the load-admittance
baseline magnitude, leaving the optimiser cost surface near-
degenerate along a curve in $(\alpha, R_x)$ space.  The closure
path is the multi-bin Taylor-Fourier estimator
(K06 PASS at $55.94\,\%$ bias improvement) on the multi-port
observation surface; the optimiser's forward-model rewire to
consume multi-bin TFT phasors on the $3 \times 3$ sending-end
admittance matrix is a follow-on commit that is not in WP5.5
scope.

\subsection{Reproducibility statement}

All code, surrogate datasets, random seeds, the Phase 4
Table~3-bis benchmark, the Phase 5 HIL pilot results, and the
v1.0 DTaaS Protection-Validation Docker image are released
under MIT licence in the public repository at
\url{https://github.com/SAMBPS-DTaaS/HIF-TF-Locator}.  Specific
reproducibility links:

\begin{itemize}
\item Public GitHub repo (HEAD onward; D0 + tags):
      \url{https://github.com/SAMBPS-DTaaS/HIF-TF-Locator}.
\item Zenodo DOIs at the release tags D2 (\texttt{v0.4.0-phase2})
      and D4 (\texttt{v0.6.0-phase4}); v1.0 DOI minted at the
      v1.0.0-phase5 tag (this paper).
\item Public CNRS IEEE-34 HIF dataset:
      \href{https://doi.org/10.57745/KRYCYY}{DOI 10.57745/KRYCYY}~\cite{CNRS_2024_IEEE34}.
\item Wang-2020 PSCAD reference open-source repo:
      \url{https://github.com/MingjieWei/PSCAD-FILE-DISTC-HIAF-Model}.
\item Docker image (v1.0): \texttt{sambps/protection-validation:v1.0}.
\end{itemize}

The full benchmark is re-runnable via
\begin{verbatim}
  python run_faultloc_phase4_benchmark.py
  python run_faultloc_phase5_hil_scenarios.py
  python -m pytest tests/ hil/ dtaas/
\end{verbatim}
which produces the per-(method, dataset) Table~3-bis CSV +
diagnostic figure + HIL scenario CSV + diagnostic plots and
runs the full $222$-test gate.

% =============================================================================
\section{Conclusion and Future Work}\label{sec:conclusion}
% =============================================================================

We presented a complete framework for single-ended joint
estimation of HIF location and arc resistance on distribution
feeders, built end-to-end across five project phases (P1
single-line baseline, P2 distributed-parameter forward model,
P3 three-phase + multi-port + IEEE feeders, P4 field-grade +
competitor benchmark, P5 HIL + productisation).  The headline
contributions are:
\begin{enumerate}
\item A closed-form three-phase distributed-parameter forward
      model with $4$-OOM tighter agreement to a 50-section
      reference than the WP3.1 brief target.
\item A multi-port CRLB on the $3 \times 3$ sending-end
      admittance matrix that is $9\times$ over-determined and
      whose proper-complex-Gaussian-ratio + joint dual-channel
      forms agree to machine precision.
\item A reproducible head-to-head benchmark across five
      candidate methods + three independent arc classes + five
      field-grade impairment classes.
\item An IEC 61850-9-2 SV + GOOSE round-trip pipeline with a
      verified $< 100$\,ms latency budget on every dev-box-mock
      scenario.
\item A v1.0 DTaaS module (REST API + CLI + scenario-engine
      plugin + UI widget + Docker) productising the framework
      for SAMBPS DTaaS Protection-Validation deployment.
\end{enumerate}

\textbf{Open work.}  The single-bin DFT identifiability floor
on the load-dominated IEEE 34 sub-sample is the binding
structural limit on K07-Phase-4 + K10 acceptance.  The
multi-bin TFT phasor + multi-port FIM closure (the optimiser's
forward-model rewire) is the immediate follow-on.  The live
HIL campaign at one of the three partner sites (NUS-GEMS,
NTU-CTSP, or IITM with Typhoon HIL fallback) closes K10
institutional signoff and the K09 hardware-side end-to-end
measurement.  The HVDC follow-on track via Amprion
extends the framework to multi-terminal MMC-HVDC fault
classification at Phase 6/7.

% =============================================================================
\bibliographystyle{IEEEtran}
\bibliography{references}
% =============================================================================

\end{document}
